% defined-in: apostol (page 208)
% entity-id: prop-rolles-theorem
% entity-type: prop
% macro: \RollesTheorem
% args: 1
% notation: #1

\hypertarget{prop-rolles-theorem}{}
\begin{theorem}[Rolle's Theorem]
\[
\TerForall \Function{f} \colon \ClosedInterval{[a, b]} \TermTo \RealNumbers
\left( \Continuous{f} \text{ on } \ClosedInterval{[a, b]} \right) \And
\left( \Function{f} \text{ on } \OpenInterval{(a, b)} \right) \And
\left( f(a) = f(b) \right)
\TermImplication \TermExists c \in \OpenInterval{(a, b)} \left( \Function{f}'(c) = 0 \right)
\]
\end{theorem}

\begin{proof}[RU]
Пусть $\Function{f}$ непрерывна на $\ClosedInterval{[a, b]}$. Согласно \WeierstrassExtremeValue{теореме Вейерштрасса}, $\Function{f}$ достигает своих экстремумов на $\ClosedInterval{[a, b]}$. Если $\Function{f}$ постоянна, то $\DerivativeAtPoint{f'}$ тождественно равна 0. Если $\Function{f}$ не постоянна, то хотя бы один из экстремумов достигается в некоторой точке $c \in \OpenInterval{(a, b)}$. В этой точке $\DerivativeAtPoint{f'(c) = 0}$ согласно лемме Ферма.
\end{proof}

\begin{proof}[EN]
Let $\Function{f}$ be continuous on $\ClosedInterval{[a, b]}$. By the \WeierstrassExtremeValue{Weierstrass Extreme Value Theorem}, $\Function{f}$ attains its extrema on $\ClosedInterval{[a, b]}$. If $\Function{f}$ is constant, then $\DerivativeAtPoint{f'}$ is identically 0. If $\Function{f}$ is not constant, at least one extremum must be attained at some $c \in \OpenInterval{(a, b)}$. At this point, $\DerivativeAtPoint{f'(c) = 0}$ by Fermat's Lemma.
\end{proof}